\let\negmedspace\undefined
\let\negthickspace\undefined
\documentclass[journal]{IEEEtran}
\usepackage[a5paper, margin=10mm, onecolumn]{geometry}
\usepackage{lmodern} % Ensure lmodern is loaded for pdflatex
\usepackage{tfrupee} % Include tfrupee package

\setlength{\headheight}{1cm} % Set the height of the header box
\setlength{\headsep}{0mm}     % Set the distance between the header box and the top of the text

\usepackage{gvv-book}
\usepackage{gvv}
\usepackage{cite}
\usepackage{amsmath,amssymb,amsfonts,amsthm}
\usepackage{algorithmic}
\usepackage{graphicx}
\graphicspath{{./figs/}}
\usepackage{textcomp}
\usepackage{xcolor}
\usepackage{txfonts}
\usepackage{listings}
\usepackage{enumitem}
\usepackage{mathtools}
\usepackage{gensymb}
\usepackage{comment}
\usepackage[breaklinks=true]{hyperref}
\usepackage{tkz-euclide}
\usepackage{listings}
\usepackage{gvv}
\def\inputGnumericTable{}

\usepackage[latin1]{inputenc}
\usepackage{color}
\usepackage{array}

\usepackage{longtable}
\usepackage{calc}
\usepackage{multirow}

\usepackage{hhline}
\usepackage{ifthen}

\usepackage{lscape}
\usepackage{circuitikz}

\begin{document}
	
	\bibliographystyle{IEEEtran}
	\vspace{3cm}
	
	\title{12.159}
	\author{EE25BTECH11042 - Nipun Dasari}
	\maketitle
	
	\renewcommand{\thefigure}{\theenumi}
	\renewcommand{\thetable}{\theenumi}
	\setlength{\intextsep}{10pt}
	
	\numberwithin{equation}{enumi}
	\numberwithin{figure}{enumi}
	\renewcommand{\thetable}{\theenumi}
	
	\textbf{Question}:\\
	The determinant
	\begin{align}
		\mydet{1+b & b & 1 \\ b & 1+b & 1 \\ 1 & 2b & 1}
	\end{align}
	evaluates to
	\begin{enumerate}[label=\alph*)]
		\item $0$
		\item $2b(b-1)$
		\item $2(1-b)(1+2b)$
		\item $3b(1+b)$
	\end{enumerate}
	
	\solution
	
	Let the determinant of the given matrix be $\Delta$. We will use row operations to simplify the matrix. Adding a multiple of one row to another does not change the determinant's value. This method is valid for $b \neq -1$.
	
	\begin{align}
		\Delta = \mydet{1+b & b & 1 \\ b & 1+b & 1 \\ 1 & 2b & 1}
	\end{align}
	
	First, we create zeros in the first column of the second and third rows using the first row as the pivot.
	
	Apply the operation $R_2 \to R_2 - \frac{b}{1+b}R_1$:
	\begin{align}
		\Delta &= \mydet{1+b & b & 1 \\ b - \frac{b}{1+b}(1+b) & (1+b) - \frac{b}{1+b}(b) & 1 - \frac{b}{1+b}(1) \\ 1 & 2b & 1} \\
		&= \mydet{1+b & b & 1 \\ 0 & \frac{(1+b)^2 - b^2}{1+b} & \frac{1+b-b}{1+b} \\ 1 & 2b & 1} \\
		&= \mydet{1+b & b & 1 \\ 0 & \frac{1+2b}{1+b} & \frac{1}{1+b} \\ 1 & 2b & 1}
	\end{align}
	
	Next, apply the operation $R_3 \to R_3 - \frac{1}{1+b}R_1$:
	\begin{align}
		\Delta &= \mydet{1+b & b & 1 \\ 0 & \frac{1+2b}{1+b} & \frac{1}{1+b} \\ 1 - \frac{1}{1+b}(1+b) & 2b - \frac{1}{1+b}(b) & 1 - \frac{1}{1+b}(1)} \\
		&= \mydet{1+b & b & 1 \\ 0 & \frac{1+2b}{1+b} & \frac{1}{1+b} \\ 0 & \frac{2b(1+b)-b}{1+b} & \frac{1+b-1}{1+b}} \\
		&= \mydet{1+b & b & 1 \\ 0 & \frac{1+2b}{1+b} & \frac{1}{1+b} \\ 0 & \frac{b+2b^2}{1+b} & \frac{b}{1+b}}
	\end{align}
	We can see that the third row is $b$ times the second row ($R_3 = b \cdot R_2$). This indicates linear dependence. To make this explicit, we perform one final operation, $R_3 \to R_3 - bR_2$.
	\begin{align}
		\Delta &= \mydet{1+b & b & 1 \\ 0 & \frac{1+2b}{1+b} & \frac{1}{1+b} \\ 0 & \frac{b(1+2b)}{1+b} - b\left(\frac{1+2b}{1+b}\right) & \frac{b}{1+b} - b\left(\frac{1}{1+b}\right)} \\
		&= \mydet{1+b & b & 1 \\ 0 & \frac{1+2b}{1+b} & \frac{1}{1+b} \\ 0 & 0 & 0}
	\end{align}
	Since the matrix has a row consisting entirely of zeros, its determinant is 0.
	\begin{align}
		\Delta = 0
	\end{align}
	The correct option is (a).
	
\end{document}